\documentclass[11pt,a4paper]{article}
\usepackage[margin=1in]{geometry}
\usepackage{amsmath,amssymb,amsfonts}
\usepackage{booktabs}
\usepackage{hyperref}
\usepackage{tikz}
\usepackage[official]{eurosym}
\usetikzlibrary{shapes.geometric, arrows.meta, positioning}

% Standardized TikZ configuration
\tikzset{
    myblock/.style={
        rectangle, 
        draw, 
        fill=blue!10, 
        text width=12cm, 
        align=center, 
        rounded corners, 
        minimum height=1.2cm,
        thick
    },
    myline/.style={
        draw, 
        thick, 
        ->, 
        >=Stealth
    }
}

\def\myperiod{.}

\def\halwebref#1{\href{https://hal.inria.fr/hal-#1}%
     {{https://hal.inria.fr/hal-#1}}}

\def\webref#1{\href{#1}%
     {{#1}}}

\newcommand{\copubsitemref}[4]{
  \item#1: 
  {\def\bfdefault{bx}\bf\emph{#2\myperiod}} #3.
  \\{\small\tt\webref{#4}}
}

\newcommand{\copubsitemhalref}[4]{
  \item#1: 
  {\def\bfdefault{bx}\bf\emph{#2\myperiod}} #3.
  \\{\small\tt\halwebref{#4}}
}
\newcommand{\copubsitemfullref}[5]{
  \item#1: 
  {\def\bfdefault{bx}\bf\emph{#2\myperiod}} #3.
  \\{\small\tt\webref{#4}}
  \\{\small\tt\halwebref{#5}}
}

\newcommand{\TODO}[1]{{\color{red}TODO: #1}}

\usepackage{marginnote}
\newcommand{\matteo}[1]{{\color{magenta}#1}}
\newcommand{\mmatteo}[1]{\marginpar{\small{\color{magenta}M: #1}}}

\newcommand{\memb}[4]{\item #1, \emph{#2}, {\tt #3}\\ #4}
\newcommand{\fmemb}[4]{\item #1, \emph{#2}, {\tt #3}}
\newcommand{\sacaddr}{1 rue Honor\'{e} d'Estienne d'Orves,
%B�timent Alan Turin,
%Campus de l'\'{e}cole Polytechnique,
Inria \& \'{E}cole Polytechnique,
91120 Palaiseau,
France}
\newcommand{\sacaddrb}{Inria Saclay, 1 rue Honor\'{e} d'Estienne d'Orves,
B\^{a}timent Alan Turin,
Campus de l'\'{E}cole Polytechnique,
%Inria \& \'{e}cole Polytechnique,
91120 Palaiseau,
France}

\title{\textbf{Project Proposal: Learning-Guided Semantics for Non-Classical Reasoning (LGS-NCR)}}
\author{\textbf{Research Track:} Artificial Intelligence}
\date{\vspace{-5ex}}

\begin{document}

\maketitle

\begin{abstract}
Modern automated reasoning systems rely heavily on a simple but powerful assumption: the truth of a complex statement can always be determined from the truth of its parts. This principle underpins classical logic and has made possible highly successful technologies such as SAT solvers, which are now used to verify hardware and software, optimise industrial processes, and solve challenging combinatorial problems.
%
However, many forms of reasoning that arise in mathematics, computer science, philosophy, and artificial intelligence do not obey this principle. Modal logics, intuitionistic logics, and logics that tolerate inconsistent information require richer semantic models, making automated reasoning dramatically more difficult. Existing methods often suffer from an explosion in the number of possibilities that must be explored, limiting their scalability to large and complex problems.
%
This project builds on Restricted Non-deterministic Matrices (RNmatrices), a powerful semantic framework that has recently enabled state-of-the-art automated theorem provers for several important families of non-classical logics. These provers can determine whether a statement is valid and automatically construct counterexamples when it is not. Despite these advances, the underlying search procedures remain computationally expensive.
%
The project proposes a fundamentally new approach by combining rigorous logical semantics with modern machine learning. Rather than replacing formal reasoning, machine learning models will learn to guide it, predicting the most promising reasoning paths while an underlying theorem prover guarantees that every result remains mathematically correct. This hybrid architecture combines the adaptability of AI with the reliability of formal verification.
%
If successful, the project will establish a new paradigm for automated reasoning in non-classical logics. Beyond advancing the foundations of logic, it will provide scalable reasoning technologies for large research programmes in formal verification, trustworthy AI, knowledge representation, and automated mathematics. More broadly, it will demonstrate how machine learning and symbolic reasoning can work together to solve problems that neither approach can tackle effectively on its own, contributing to one of the central challenges in artificial intelligence.
%The principle of truth-functionality, central to classical and many-valued logics, states that the truth-value of a complex formula is uniquely determined by the truth-values of its subformulas. This makes truth tables both easy to understand and implement, enabling SAT solvers to crack remarkably hard problems. But this tidy picture cannot be obtained when crossing the boundary of classicality, reaching e.g., modal logics. Non-deterministic matrices (Nmatrices) generalise ordinary matrices by allowing connectives to return a set of possible truth values rather than a unique one, thereby increasing expressive power. Restricted non-deterministic matrices (RNmatrices) further refine this framework by imposing constraints on the rows of Nmatrices, filtering out ``unsound'' rows and retaining only ``valid'' ones. This yields a more expressive semantic framework that has been successfully used to provide sound and complete semantics for a wide range of logics, including paraconsistent, intuitionistic, and modal logics. Recently,  automated theorem provers (ATP) based on RNmatrices were developed by encoding their semantics as Satisfiability Modulo Theories (SMT) problems. The resulting provers decide validity and construct countermodels, achieving state-of-the-art performance for paraconsistent logics and competitive results for intuitionistic and modal logics. While this framework yields robust decision procedures, it still suffers from a combinatorial explosion due to the vast space of admissible rows, non-deterministic valuation choices, and structural elimination criteria. This project proposes a novel hybrid paradigm that embeds Machine Learning (ML) techniques directly into the semantics-driven proof search of RNmatrices. The project aims to dramatically prune the search space, by training graph neural networks (GNNs) and transformer models to predict viable row constraints, prioritize candidate valuations, and anticipate minimal countermodels. Crucially, the underlying SMT-solver back-end preserves absolute soundness and completeness, uniting data-driven adaptive heuristics with formal proof guarantees.
\end{abstract}

\subsection*{1. Presentation of the teams}

\subsubsection*{a. French Team}

The French Team consists of members of the {\sc Partout} project which is joint between Inria Saclay and LIX (UMR 7161) at \'{E}cole Polytechnique. This team brings extensive expertise in all areas of logic, formal verification, and proof theory, which is crucial for the success of the project.

\begin{itemize}
\fmemb{PI: Kaustuv Chaudhuri}{Charg\'{e} de recherche}{kaustuv@chaudhuri.info}{\sacaddr}

\fmemb{Dale Miller}{Directeur de recherche}{dale.miller@inria.fr}{\sacaddr}

\fmemb{Lutz Stra\ss burger}{Directeur de recherche}{lutz@lix.polytechnique.fr}{\sacaddr} 

\fmemb{Beniamino Accattoli}{Charg\'{e} de recherche}{beniamino.accattoli@inria.fr}{\sacaddr}

\fmemb{Victoria Barrett}{Postdoc}{victoria.barrett@inria.fr}{\sacaddr}

\fmemb{Arunava Gantait}{PhD student}{Arunava.Gantait@lix.polytechnique.fr}{\sacaddr}

\end{itemize}
The postal address for all team members is:\\[.5ex] \sacaddrb.


\subsubsection*{b. British Team}

The British team is located at University College London, and it comprises experts in logic and verification, as well as machine learning and proof theory. 
\begin{itemize}

\fmemb{PI: Elaine Pimentel}{Professor of logic and Computation}{e.pimentel@ucl.ac.uk}{66-72 Gower St., University College London,
WC1E 6EA London, UK}

\fmemb{David Pym}{Professor of Information, Logic, and Security}{d.pym@ucl.ac.uk}{66-72 Gower St., University College London,
WC1E 6EA London, UK}

\fmemb{Victor Barroso Nascimento}{Postdoc}{victorluisbn@gmail.com}{66-72 Gower St., University College London,
WC1E 6EA London, UK}

\fmemb{Katya Piotrovskaya}{PhD student}{kate.piotrovskaya.21@ucl.ac.uk}{66-72 Gower St., University College London,
WC1E 6EA London, UK}

\fmemb{Maria Os\'{o}rio}{PhD student}{ucacmo0@ac.ucl.uk}{66-72 Gower St., University College London,
WC1E 6EA London, UK}

\fmemb{Joaquim Torres Waddington}{PhD student}{joaquim.waddington.24@ucl.ac.uk}{66-72 Gower St., University College London,
WC1E 6EA London, UK}

\fmemb{Arghya Chackrabarty}{PhD student}{arghya.chakrabarty.20@ucl.ac.uk}{66-72 Gower St., University College London,
WC1E 6EA London, UK}

\end{itemize}
%
The postal address for all members of the British team is:\\[.75ex]
66-72 Gower St., University College London,
WC1E 6EA London, UK



%\TODO{make smooth; maybe move after the project}

\subsection*{2. The Research Project}

\paragraph{State of the Art:}

%\subsection{The Challenge of Truth-Functionality in Non-Classical Logics}
In classical propositional logic, truth-functionality ensures that the truth value of a complex formula is uniquely determined by its subformulas. This structural rigidity underpins modern SAT and SMT technologies. However, when moving beyond classical bounds into constructive or modal realms, finite truth-functionality collapses. G\"{o}del proved in 1932 that no finite-valued deterministic matrix can faithfully characterize propositional intuitionistic logic ($IL_p$). Similarly, Dugundji demonstrated in 1940 that finite truth-functionality cannot preserve modal interpretations of ``qualified truth'' across the modal cube.

To overcome these constraints without resorting to infinite-valued structures, Avron and Lev introduced Non-Deterministic Matrices (Nmatrices), substituting truth functions with multifunctions that map tuples of values to a set of possible values.

%\subsection{From Nmatrices to RNmatrices}
While Nmatrices capture many non-classical behaviors, they are occasionally too permissive. For instance, blind non-determinism fails to model the strict proof-theoretic evolution of intuitionistic logic or the rule of necessitation in modal logic, generating invalid matrix rows that break soundness. Restricted Non-Deterministic Matrices (RNmatrices) resolve this by enforcing global first-order or modal constraints over the set of allowed valuations, filtering out non-deterministic choices that do not align with the target deductive system\footnote{R. R. Leme, C. Olarte, E. Pimentel and M. E. Coniglio, {\em The Modal Cube Revisited: Semantics without Worlds}, TABLEAUX 2025.}. 

%Table~\ref{tab:logics} summarizes key semantic components across target formalisms.
%
%\begin{table}[htbp]
%\centering
%\caption{Core Languages and Truth-Value Dimensions in RNmatrix Semantics}
%\label{tab:logics}
%\begin{tabular}{llll}
%\toprule
%\textbf{Logic} & \textbf{Core Language Elements} & \textbf{Truth Values ($\mathcal{V}$)} & \textbf{Designated Values ($\mathcal{D}$)} \\ \midrule
%\textbf{Modal $S4$} & $p, \perp, \top, \neg, \Box, \rightarrow, \wedge, \vee$ & $\{0, 1, 2\}$ & $\{1, 2\}$ \\
%\textbf{Intuitionistic $IL_p$} & $p, \perp, \top, \neg, \rightarrow, \wedge, \vee$ & $\{T, F\}$ & $\{T\}$ \\
%\textbf{Paraconsistent $C_n$} & $p, \neg, \rightarrow, \wedge, \vee$ & $\{T_n, t_0^n, \dots, t_{n-1}^n, F_n\}$ & $\mathcal{V}_{C_n} \setminus \{F_n\}$ \\ \bottomrule
%\end{tabular}
%\end{table}

%\subsection{The TRiNity SMT Framework}
Recent progress materialized in the automated theorem prover TRiNity\footnote{R. R. Leme, C. Olarte, E. Pimentel, {\em Efficient Decision Procedures for RNmatrix Semantics}, LSFA 2026.}, which translates validity problems into a multi-sorted SMT signature containing sorts for truth values (\texttt{Val}), rows (\texttt{Row}), and formula columns (\texttt{Col}). A formula $\alpha$ is valid if and only if an arbitrary row $r_0$ cannot falsify it under the structural matrix constraints and the logic-specific restriction criteria.

%\section{Problem Statement \& Motivation}

\paragraph{Research programme:}
Although the SMT-encoding methodology outperforms standard tableau-based systems for logics like paraconsistent da Costa's $C_n$ hierarchy, it encounters severe bottlenecks when scaled up:
\begin{enumerate}
    \item \textbf{Combinatorial Row Explosion:} In logics requiring nested or backward-looking existence criteria (such as $IL_p$ and modal $S4$), verifying a valuation necessitates confirming the existence of companion or successor valuations. In SMT-LIB syntax, this translates to heavily quantified formulas ($\forall r . \exists w$) over uninterpreted matrix functions (\texttt{mat}).
    \item \textbf{Blind Valuation Guessing:} SMT solvers rely on generic, logic-agnostic heuristic engines (e.g., CDCL, internal simplex variants, or E-matching for quantifiers). They do not exploit the highly structured semantic layouts inherent to matrix representations.
    \item \textbf{Countermodel Inflation:} When a formula is invalid, the SMT solver searches for a countermodel. Without domain guidance, the constructed countermodels are frequently bloated and suboptimal, obscuring the core structural reasons for invalidity.
\end{enumerate}

%\section{Proposed Methodology: The Hybrid AI-SMT Framework}
We propose a hybrid framework that injects intelligent, data-driven heuristics into the RNmatrix execution pathway while preserving the exact completeness and correctness of the SMT layer. 

%The schema below illustrates the system architecture.
%\begin{center}
%\begin{tikzpicture}[node distance=1.8cm, auto]
%    % Nodes (explicitly utilizing the corrected 'myblock' style)
%    \node [myblock] (input) {\textbf{Input Non-Classical Formula}};
%    \node [myblock, below=of input] (ai) {\textbf{Deep Neural Heuristic Engine (GNN / Transformer)} \\ - Predicts sound row valuation patterns \\ - Prioritizes promising proof paths \& countermodels};
%    \node [myblock, below=of ai] (trinity) {\textbf{TRiNity Matrix-to-SMT Encoder} \\ Emits multi-sorted SMT-LIB files annotated with AI hints};
%    \node [myblock, below=of trinity] (smt) {\textbf{SMT Back-End Solver (Z3 / CVC5)} \\ Enforces formal constraints \& absolute correctness};
%    
%    % Paths (explicitly utilizing the corrected 'myline' style)
%    \path [myline] (input) -- (ai);
%    \path [myline] (ai) -- node [right, font=\small, text width=6cm, align=left] {Neural Guidance \\ (Weights \& Hints)} (trinity);
%    \path [myline] (trinity) -- (smt);
%\end{tikzpicture}
%\end{center}
%\subsection{Neural Structure Encoding}
Logical formulas will be parsed as Heterogeneous Graph structures where subformula dependencies represent edges, and operators ($\Box, \neg, \rightarrow$) represent node types. Graph Neural Networks (GNNs) will learn low-dimensional vector representations of these formulas. In parallel, a Transformer model will sequence the subformula tokens to learn global contextual boundaries.

%\subsection{Semantic Valuation \& Row Prioritization}
Instead of forcing the SMT solver to blindly evaluate the uninterpreted function \texttt{mat(Row, Col)}, the trained model outputs a probability distribution over the matrix sort space:
\begin{itemize}
    \item \textbf{Row Selection:} Predicts which rows in a multi-row evaluation (e.g., Kripke-like accessibility pathways in $IL_p$ or $S4$) are vital to trigger a contradiction or satisfy an assertion.
    \item \textbf{Valuation Biasing:} Injects SMT \texttt{set-info} or branch-fuzzing weights that bias the solver toward evaluating designated or non-designated values based on structural formula features.
\end{itemize}

%\subsection{Predictive Countermodel Generation}
For invalid formulas, the neural model will be optimized using a customized loss function to output a partial matrix assignment that minimizes the required valuation cardinality. The SMT solver uses this partial matrix as a ``warm-start'' hint, converting a complex synthesis task into an efficient checking task.

%\subsection{Soundness and Completeness Preservation}
Crucially, AI predictions are treated solely as soft constraints (search biases, ordering heuristics, or initial assignments). If the neural model makes an incorrect inference, the underlying SMT core ensures it is safely bypassed, guaranteeing that the prover remains fully sound and complete.

%\section{Evaluation Plan \& Benchmarking}
The performance of the learning-guided framework will be measured quantitatively across three execution domains:
\begin{enumerate}
    \item \textbf{The Paraconsistent Hierarchy ($C_n$):} Benchmarking on deep nested depth contradictions ($c_k(\alpha)$). We will assess if the AI can predict the exact point where the ``classical core'' triggers explosion ($k \ge n$), thereby skipping unnecessary matrix searches.
    \item \textbf{Intuitionistic Logic ($IL_p$):} Evaluation on the intuitionistic ILTP library. The target metric is the time reduction in validating the localized existence criterion for implication and negation ($\forall v \exists w$).
    \item \textbf{The Modal Cube ($S4$):} Evaluating performance on modal formulas that test iterated necessity ($\Box\Box\alpha \leftrightarrow \Box\alpha$). The framework will be optimized to efficiently navigate Gr\"{a}tz's partial level valuations.
\end{enumerate}

Success will be defined as an increase in the number of solved problems within a standard 300-second timeout window, combined with a reduction in back-tracking steps inside the SMT engine compared to the baseline TRiNity prover.

%\section{Expected Impact \& Future Work}
This project introduces a scalable approach to automated non-classical reasoning by combining data-driven predictive power with the formal robustness of SMT-encoded RNmatrices. Beyond enhancing proving speeds, the framework will automate the discovery of minimal semantic countermodels, rendering them highly interpretable for computer scientists and logicians. Future iterations will seek to expand this learning-guided semantic paradigm to non-propositional extensions, first-order modal systems, and dynamic epistemic frameworks.

\paragraph*{Historic Context of the Collaboration:}
The longstanding collaboration between the PARTOUT group and the PPLV team at UCL dates back to the early 2000s, when Dale Miller, one of the founders of Parsifal (the predecessor of PARTOUT) mentored Elaine Pimentel. At that time, Pimentel was an Assistant/Associate Professor at UFMG and later became a Full Professor at UFRN in Brazil. She joined UCL in 2022, where she currently holds the position of Professor of Logic and Computation.

Since those early years, the collaboration between Parsifal/PARTOUT and Pimentel has remained active and productive. Between 2002 and 2012, Lutz Stra\ss burger, Alexis Saurin, and Dale Miller, members of the Parsifal group, made five research visits to UFMG and UFRN, while Brazilian researchers regularly visited France, typically for one- to two-month stays each year.

In 2018, Pimentel spent part of her sabbatical year at Inria Saclay, further strengthening the academic ties between the teams. During this period, her collaboration with Chaudhuri began, culminating in a particularly significant publication\footnote{K. Chaudhuri, J. Despeyroux, C. Olarte and E. Pimentel, {\em Hybrid Linear Logic, revisited}  MSCS Volume 29 Issue 8, pp. 1151-1176, 2019.}, a joint paper written in honor of Dale Millers 60th birthday.

This enduring collaboration is evidenced by several joint publications and continues to thrive through the joint organization of major international seminars, including the BANFF (2025) and Dagstuhl (2024) meetings on \emph{Proof Representations: From Theory to Applications} (co-organized with Lutz Stra\ss burger).

The depth of this partnership is further underscored by active participation in several PhD defense committees at Inria and \'{E}cole Polytechnique, including those of E. Grienenberger (2025), J. Wu (2024), M. Morales (2023), S. Marin, and V. Nigam.


%In addition, active participation as jury member or rapporteur in several Parsifal/PARTOUT PhD defences highlights the depth and continuity of these collaborations:
%%
%\begin{itemize}
%\item \'{E}milie Grienenberger -- Inria Paris-Saclay (02/2025)
%\item Jui-Hsuan (Ray) Wu -- \'{E}cole Polytechnique (12/2024)
%\item Marianela Morales -- \'{E}cole Polytechnique (01/2023) -- reviewer
%\item Sonia Marin -- \'{E}cole Polytechnique (01/2023, 01/2018)
%\item Vivek Nigam -- \'{E}cole Polytechnique (01/2023, 09/2009)
%\end{itemize}

Collectively, these activities demonstrate the strength, longevity, and productivity of the research partnership between PARTOUT and PPLV.


\paragraph{Plans for continued collaboration and funding prospects after the project:}

%, notably a potential Horizon Europe application}

The project is intended to create a critical mass necessary to launch a larger program under Horizon Europe. We plan to organise workshops and invite external experts to eventually enlarge the consortium. A successful outcome of this project will allow us to attract further funding for the research programme outlined above.
This initiative aims at building a strong foundation and demonstrating the potential impact,  and to secure further funding and support for extensive research and development efforts within the Horizon Europe framework.


\paragraph{Long term outcomes of the project:} %%, including potential societal and industrial outcomes}
The successful execution of this project is expected to yield substantial, long-term impacts across both the automated reasoning and formal verification communities:

\begin{itemize}
    \item \textbf{A Standardized API for AI-Guided Semantic Provers:} The project will result in an open-source, modular extension to the TRiNity engine. By standardizing the interface through which neural models pass valuation and row heuristics to SMT back-ends, we establish an extensible infrastructure that allows other researchers to swap in different neural architectures (e.g., Graph Transformers vs. Reinforcement Learning agents).
    \item \textbf{Universal Benchmark Libraries for Non-Classical Logics:} To train the models, the project will generate a massive dataset of non-classical formulas across the $C_n$ hierarchy, intuitionistic logic, and modal systems, mapping formulas to their corresponding RNmatrix search paths and minimal countermodels. This will serve as a permanent benchmark library for future AI-driven ATP systems.
    \item \textbf{Accelerated Formal Verification for Specialized Systems:} Non-classical logics are critical for verifying distributed networks (modal logic), cryptographic protocols (intuitionistic logic), and database systems handling contradictory data (paraconsistent logic). Providing an order-of-magnitude speedup in these domains will make practical formal verification viable for complex, safety-critical industrial applications that classical SAT/SMT tools cannot naturally handle.
\end{itemize}

\paragraph*{Risk Analysis \& Mitigation Strategies}
Given the interdisciplinary nature of fusing deep learning models with formal symbolic reasoning, several technical and structural risks have been identified alongside concrete mitigation pathways.

%\begin{table}[htbp]
%\centering
%\caption{Risk Assessment and Mitigation Matrix}
%\label{tab:risks}
%\begin{tabular}{p{3.5cm}llp{6.5cm}}
%\toprule
%\textbf{Identified Risk} & \textbf{Probability} & \textbf{Impact} & \textbf{Mitigation Strategy} \\ \midrule
%\textbf{Data Sparsity for Training} & Medium & High & Use synthetic data generation engines by mutating existing templates and utilizing structural bootstrapping techniques. \\ \addlinespace
%\textbf{Neural Inference Latency Bottleneck} & High & Medium & Implement lightweight, distilled models (e.g., highly compressed GNNs) and restrict neural inference to complex subformulas above a specified nesting threshold. \\ \addlinespace
%\textbf{Poor Heuristic Convergence} & Medium & Low & \textbf{Architectural Safeguard:} The SMT layer guarantees absolute correctness. If neural guidance degrades, the engine falls back gracefully to standard TRiNity search heuristics. \\ \bottomrule
%\end{tabular}
%\end{table}
\begin{enumerate}
    \item \textbf{The Data Sparsity Challenge:} Unlike classical logic, standard libraries for non-classical systems (like the ILTP library) are relatively small. Training deep neural networks requires massive quantities of labeled data. 
    \\ \textit{Mitigation:} We will implement a synthetic formula generator that systematically construct valid and invalid formulas by mutating known matrix properties, thereby creating a balanced and rich training space.
    \item \textbf{Inference Overhead vs. Prover Speed:} Complex graph or transformer model passes could introduce a runtime overhead that cancels out the speedups gained by pruning the SMT search space.
    \\ \textit{Mitigation:} Models will be aggressively quantized and distilled into lightweight edge architectures. Furthermore, neural prediction will only trigger at specified structural decision points (e.g., when encountering heavily quantified row constraints or multi-level modal jumps).
    \item \textbf{Heuristic Misdirection (Hallucination):} The neural model might predict faulty row selections or misleading valuation biases, guiding the prover into highly inefficient search paths.
    \\ \textit{Mitigation:} Because soundness and completeness are fully decoupled from the neural layer, bad predictions never compromise proof integrity. If a search path exceeds a localized timeout, the prover automatically dampens the neural weights and defaults to standard SMT backtracking.
\end{enumerate}



\section*{Annexes}

\subsection*{Calendar for the project}

\begin{itemize}
    \item Start date: 1 January 2027
    \item Duration 16 months
    \item Main milestones: We expect each of the four work packages to take about 3-4 months, giving us room to adjust in case of unforeseen difficulties. The last two month are reserved for writing up our result in form of publications in prestigious conferences and journals.
    \begin{itemize}
        \item  April 2027: Build the Deep Neural Heuristic Engine (GNN / Transformer)
        \item July 2027: Propose the TRiNity Matrix-to-SMT Encoder
        \item November 2027: Implement the SMT Back-End Solver (Z3 / CVC5)
        \item 2027: Benchmark the results
    \end{itemize}
\end{itemize}


\subsection*{Financial details}

\subsubsection*{a. Planned budget}

\par
%Observe that: ``The funds available in this call are provided by the French government and will be allocated to French researchers. Although this funding is limited to French researchers, only applications jointly submitted by French and British PIs to the French Steering Committee are declared eligible. British researchers will have to apply to the British call, which will be launched separately from the French one.''

%Hence the budget should be only for the French side.

\smallskip

%FRANCE

\begin{tabular}{|l|r|l|}
\hline
Type of expense & Amount (in \euro) & Description \\
\hline
Travel & 22 800  & mutual visits and conference attendance\\
Equipment & 3 000 & 2 laptops \\
Essential consumable materials & 200 & \\
Documentation &  - & \\
Publication & 2 000 & \\
%Field work (if relevant) & & \\
Research grants & - & \\ %%two one-month research visits abroad\\
Other & 11 000 & workshop organisation\\
\hline
\bf Total & \bf 39 000 & \\
\hline
\end{tabular}


\bigskip

\subsubsection*{b. Justification of expenses}

\subsection{Justification of Expenses}
The proposed budget is directly tied to the core objectives of the project, focusing on maximizing the collaborative expertise of the two teams, enabling intensive neural network training, and fostering wider research consortiums. 

\begin{itemize}
    \item \textbf{Travel and Subsistence (\euro 22,800):} Since this is a joint initiative between PARTOUT (France) and PPLV (UK), face-to-face collaboration is paramount. The bulk of the budget is allocated to cross-border research exchanges, specifically prioritizing PhD students and junior researchers. 
    \begin{itemize}
        \item \textit{Mutual Bilateral Visits (\euro 17,800):} Covers 10 distinct one-week research visits (5 from France to the UK, and 5 from the UK to France). Each visit is budgeted at \euro 700, encompassing \euro 200 for return train travel and \euro 500 for local accommodation and subsistence (\euro 100/day for 5 days). It also covers 4 one-month visits from PhD students and postdocs, each visit is budgeted at \euro 2,700: \euro 200 for return train travel and \euro 2,500 for local accommodation and subsistence.
        \item \textit{Conference Attendance (\euro 5,000):} Allocates \euro 2,500 each for two PhD students to present research results at major international peer-reviewed conferences. This includes registration fees (\euro 800), transportation (\euro 700), and accommodation/subsistence (\euro 1,000).
    \end{itemize}

    \item \textbf{Equipment (\euro 3,000):} To support the computational overhead of training Graph Neural Networks (GNNs) and benchmarking Large Language Models (LLMs) locally, the project requires 2 high-performance laptops equipped with dedicated GPU acceleration capabilities, budgeted at approximately \euro 1,500 each.

    \item \textbf{Consumables and Materials (\euro 200):} Covers essential day-to-day coordination materials for workspace collaboration (whiteboards, markers, and paper) as well as dedicated hardware (microphones and high-definition cameras) to ensure seamless, high-quality video-conferencing during hybrid team meetings.

    \item \textbf{Publication Costs (\euro 2,000):} While the teams strongly advocate for green open-access and host all primary preprints freely on HAL, this fund is strictly reserved to cover mandatory Open Access or processing fees levied by prestigious international conferences in automated reasoning.

    \item \textbf{Workshop Organisation (\euro 11,000):} To establish the critical mass required for future Horizon Europe bids, the project will host two milestone workshops--one in France and one in the United Kingdom. 
    \begin{itemize}
        \item \textit{External Expert Travel (\euro 8,400):} Covers the travel and lodging for two invited international experts per workshop (\euro 1,500 airfare + \euro 600 for a 3-night hotel stay per expert).
        \item \textit{Catering and Refreshments (\euro 2,600):} Provides lunch (\euro 450/day) and two coffee breaks (\euro 90/break) across the 2-day duration of each workshop for an estimated 30 participants. Institutional lecture halls are provided free of charge by the host universities.
    \end{itemize}
\end{itemize}

\subsection*{Publications of PIs}

\subsubsection*{a. Kaustuv Chaudhuri (PI of the French team)}

\begin{enumerate}

\copubsitemfullref{Kaustuv}{Title}{Venue}{doi}{HAL\#}
\end{enumerate}

\subsubsection*{a. Elaine Pimentel (PI of the British team)}

\begin{enumerate}

\copubsitemref{Elaine}{Title}{Venue}{doi}
\end{enumerate}
\end{document}